\section{Polynomial unrestricted learning of shallow targets at every depth}
\label{sec:shallow-unrestricted}

A shallow target can have a small first-layer Gaussian energy at a
polynomial base width simply by adding zero neurons. Its readout energy
is then polynomial. Distributing that readout cost over the learner's
remaining layers makes the product-norm parameter polynomial uniformly
in depth. This gives a family for which neither arbitrary width nor
arbitrary added depth requires a robustness restriction.

\begin{theorem}[Shallow targets, all depths and all admissible widths]
\label{thm:shallow-all-depths}
Suppose a Boolean target $f:\{-1,1\}^d\to\{-1,1\}$ has a one-hidden-layer
representation
\[
 z(x)=\sum_{j=1}^{C}v_j\tanh(a_j^T\bar x),\qquad
 \bar x=(1,x),\qquad f(x)z(x)\ge1
\]
with every coordinate of $a_j$ and every $v_j$ bounded in absolute value
by $R\ge1$. Let $C_L(f)$ denote its bounded-weight network complexity
at depth $L$: the least width $N$ admitting a realization with all
weights at most $N$ in absolute value and signed output greater than
$1/2$ everywhere.

For each fixed $\zeta>0$, the original unrestricted Gaussian learner has
sample complexity polynomial in
\[
 d,C,R,L,\varepsilon^{-1},\log(1/\delta),
\]
for every specified learner depth $L\ge2$ and every input distribution,
simultaneously for every width $N\ge C_L(f)$ and every
$\gamma\ge1+\zeta$. With probability at least
$1-\delta$ over a sample of that size, posterior Gibbs error is at most
$\varepsilon$ simultaneously at all these widths and all prior scales
$0<\sigma^2\le N^{-(1+\zeta)}$. The polynomial's degree may depend on
$\zeta$. The prior is conditioned only on the original global clean
margin and the training labels.
\end{theorem}

\begin{proof}
Let $A$ have rows $a_j^T$, and set
\[
 M=\left\lceil\max\{4,C,2R,d\|A\|_F^2\}\right\rceil,
 \qquad T=\max\{1,M\|v\|_2^2\}.
\]
Both are polynomial in $d,C,R$. Add zero hidden neurons to obtain width
$M$. This preserves the function, and its first-layer and readout
energies satisfy
\[
 E_{\rm pre}=\frac dM\|A\|_F^2\le1,
 \qquad c^2=M\|v\|_2^2\le T.
\]
In fact $M\|v\|_2^2\ge1$, since unit margin and $|\tanh|\le1$ imply
$\|v\|_2\ge1/\sqrt C$, but the maximum makes this unneeded.

For any $N\ge M$, balanced cloning of this shallow network preserves its
output exactly and gives
\[
 E_{\rm pre}\le2,\qquad c^2=N\|v_N\|_2^2\le2T.
\]
Explicitly, assign each of the $M$ original neurons either
$\lfloor N/M\rfloor$ or $\lceil N/M\rceil$ clones, copy its incoming
weights, and divide its readout weight by its number of clones.
The first cost increases by at most $M\lceil N/M\rceil/N\le2$,
and the readout cost by at most $N/(M\lfloor N/M\rfloor)\le2$.
For learner depth $L=2$ we already have a unit-margin witness with
\[
 K_2=2+2T,\qquad
 Q_*^2=\left(\frac{K_2+2}{2}\right)^2=(T+2)^2.
\]

For $L\ge3$, apply Lemma~\ref{lem:scalar-transmission} with
$D=4c^2$ and $m=L-1$. Replace the old readout by a layer whose every
row is $2\alpha v_N^T$. Follow it with the scalar transmission layers,
implemented by equal coordinates with dense matrix entries $b_s/N$,
and final readout entries $q/N$, where
\[
 q^2=\max\{c^2,1\}^{1/(L-1)}\le(2T)^{1/(L-1)}.
\]
The target had margin at least one, hence in particular at least
$1/2$, so the scalar construction gives margin at least $1/2$.
Double its final readout to restore margin at least one. This adds
$3q^2$ to the energy. Thus every width $N\ge M$ has a depth-$L$
unit-margin witness with
\[
 \kappa_N\le K_L:=2+(L+14)(2T)^{1/(L-1)}.
\]
Writing $q_*^2=(2T)^{1/(L-1)}\ge1$, we have
$K_L+2\le(L+18)q_*^2$. Consequently
\begin{equation}\label{eq:shallow-product-polynomial}
 Q_*^2=\left(\frac{K_L+2}{L}\right)^L
 \le e^{18}(2T)^{L/(L-1)}
 \le e^{18}(2T)^{3/2}\qquad(L\ge3).
\end{equation}
The exponent $L/(L-1)$ is at most $3/2$ for $L\ge3$.
Theorem~\ref{thm:unrestricted-above-one}, with base width $M$, therefore
provides a sample bound polynomial in the stated quantities for all
$N\ge M$.

We verify admissible smaller widths rather than silently identifying
$M$ with the target's minimum complexity. First, $C_L(f)\le M$ for all
$L\ge2$. At width $M$, the original incoming weights have magnitude
at most $R\le M/2$. In the added scalar channel,
$\alpha\le\operatorname{atanh}(1/2)<0.55$, so its first weights
$2\alpha v_j$ have magnitude at most $1.1R<M$.
The scalar construction has $b_s\le2.2q$ and final weights $2q/M$.
Since $T\le MC R^2\le M^4$ and $L\ge3$,
\[
 q\le(2T)^{1/4}\le2^{1/4}M,
 \qquad b_s/M<3\le M,\qquad 2q/M<3\le M.
\]
Thus this depth-$L$, width-$M$ circuit satisfies the bounded-weight
convention and has margin at least one. The case $L=2$ is immediate.

For $C_L(f)\le N<M$, start from a defining width-$C_L(f)$ circuit,
clone it to width $N$, and double its readout. It has margin greater
than one, all weights at most $2M$, and Gaussian energy bounded by a
polynomial in $d,L,M$. The Gaussian box estimate and
Theorem~\ref{thm:finitewidth-unrestricted} consequently give a sample
bound polynomial in $d,L,M,\varepsilon^{-1},\log(1/\delta)$,
simultaneously for all these widths and all $t\le1$.
Use failure probability $\delta/2$ for this event and $\delta/2$ for
the preceding event for $N\ge M$. Their intersection proves the claim.
\end{proof}

The complexity parameter controlling this theorem's sample size is the
size and weight bound of a \emph{shallow} representation. Its width
admissibility is the actual depth-$L$ complexity $C_L(f)$. The theorem
does not identify shallow complexity with the minimum complexity of an
arbitrary deep circuit, and therefore does not settle the unrestricted
problem for all small deep circuits. It also makes no claim at $\gamma=1$.

\begin{corollary}[Parity at arbitrary depth and admissible width]
\label{cor:parity-full-polynomial}
For every fixed $\zeta>0$, parity of any $k\ge1$ chosen coordinates in
$\{-1,1\}^d$ has unrestricted Gaussian posterior sample complexity
polynomial in $d,k,L,\varepsilon^{-1},\log(1/\delta)$, for every
$L\ge2$, every $N\ge C_L(f)$, and every $\gamma\ge1+\zeta$.
The bound holds for every input distribution and for the original
finite-variance posterior.
\end{corollary}

\begin{proof}
Let $q=(k-\sum_{i\in S}x_i)/2\in\{0,\ldots,k\}$ count negative
coordinates, and put $a=\log(16k)$. Define
\[
 \tau_j(x)=\tanh\bigl(a(q-j+1/2)\bigr),\qquad 1\le j\le k.
\]
Every argument has magnitude at least $a/2$, so
$|\tau_j-\operatorname{sign}(q-j+1/2)|\le2e^{-a}=1/(8k)$.
Writing $c_k=1$ for even $k$ and $c_k=0$ for odd $k$, the exact step
identity is
\[
 (-1)^q=c_k+\sum_{j=1}^k(-1)^j
                  \operatorname{sign}(q-j+1/2).
\]
Hence $F=c_k+\sum_j(-1)^j\tau_j$ satisfies
$|F-(-1)^q|\le1/8$. Represent the constant by one extra hidden neuron
with augmented-input weight $\operatorname{atanh}(1/2)$ and no other
incoming weights. A linear readout equal to $2F$ has signed margin
at least $7/4$. Its threshold-neuron readout weights are $2(-1)^j$,
and its constant-neuron weight is $4c_k$.
Its incoming weights have magnitude at most $a(k+1)$.
Thus $C=k+1$ and $R=\max\{4,a(k+1)\}$ are polynomial in $k$.
Apply Theorem~\ref{thm:shallow-all-depths}.
\end{proof}

\begin{corollary}[Symmetric Boolean targets]
The conclusion of Corollary~\ref{cor:parity-full-polynomial} holds for
any Boolean function of the Hamming weight of $k$ selected coordinates.
\end{corollary}

\begin{proof}
If its values are $h(0),\ldots,h(k)\in\{-1,1\}$, then
\[
 h(q)=\frac{h(0)+h(k)}2+
       \sum_{j=1}^k\frac{h(j)-h(j-1)}2
                    \operatorname{sign}(q-j+1/2).
\]
Each coefficient has magnitude at most one. Replacing the signs by
the same $\tau_j$ and doubling the output gives margin at least $7/4$,
using the same width and weight bounds as in the parity proof.
\end{proof}


\begin{corollary}[Polynomial-size DNF and CNF targets]
A Boolean function on $d$ bits represented by a DNF or CNF with $s\ge1$
clauses has unrestricted Gaussian posterior sample complexity polynomial
in $d,s,L,\varepsilon^{-1},\log(1/\delta)$, for every specified $L\ge2$,
every $N\ge C_L(f)$, and every $\gamma\ge1+\zeta$, with $\zeta>0$ fixed.
\end{corollary}

\begin{proof}
For a DNF, let $q_j(x)$ be the number of false literals in clause $j$,
and put $a=\log(16s)$ and
$\tau_j(x)=\tanh(a(1/2-q_j(x)))$.
The exact clause indicator is
$[1+\operatorname{sign}(1/2-q_j)]/2$.
Thus
\[
 H(x)=s-1+\sum_{j=1}^s\operatorname{sign}(1/2-q_j(x))
\]
is $-1$ when no clause is satisfied and at least $1$ otherwise.
Replacing the signs by $\tau_j$ changes $H$ by at most $1/8$.
Doubling the resulting expression gives global signed margin at least
$7/4$. It is a shallow tanh network with $s+1$ features: the extra
feature is the constant $1/2$, with readout weight $4(s-1)$.
After removing duplicate literals and discarding contradictory clauses,
each clause has at most $d$ literals. Each $q_j$ is affine in the bits,
so the remaining incoming and readout weights have magnitude at most
$R=\max\{4s,a(d+1)\}$. The shallow-target theorem applies.
A CNF is the negation of a DNF in negated literals, and negating the
readout preserves every bound. Constant functions are immediate.
\end{proof}
